% ============================================================
% PRN232 - Section 3: System Architecture (Vietnamese & English)
% ============================================================
\section{\bi{Kiến trúc hệ thống}{System Architecture}}
\label{sec:arch}

\subsection{\bi{Tổng quan phân tầng}{Layered Overview}}

\bi{
Hệ thống được tổ chức theo bốn tầng, mỗi tầng chỉ giao tiếp với tầng
ngay bên dưới. Tầng trình bày gồm cổng web Next.js 15 (33 route) và ứng dụng di
động Flutter (31 màn hình); tầng API là một tiến trình NestJS duy nhất với 19 module
nghiệp vụ và hơn 140 endpoint REST; tầng thời gian thực xử lý phát sóng trực tiếp (LiveStream), điều phối phòng Studio, trò chuyện, cảm xúc nổi, quyên góp Super Chat và tín
hiệu WebRTC; tầng dữ liệu gồm PostgreSQL (21 model Prisma), máy chủ phát trực tuyến
MediaMTX, dịch vụ lưu trữ đám mây Cloudinary và các API bên ngoài.
}{
The system is organised into four layers, each communicating only with
the layer immediately below it. The presentation layer comprises the
Next.js 15 web portal (33 routes) and the Flutter mobile application (31 screens); the API
layer is a single NestJS process with 19 business modules and 140+ REST endpoints; the
real-time layer handles live broadcasting (LiveStream), studio room orchestration, live chat, floating reactions, Super Chat donations and WebRTC signalling;
and the data layer consists of PostgreSQL (21 Prisma models), the MediaMTX streaming
server, Cloudinary storage and external APIs.
}

\begin{figure*}[t]
\centering
\fitfull{\figArchitecture}
\caption{\bi{Kiến trúc phân tầng tổng thể của OmniCast. Cả cổng web Next.js 15 và
ứng dụng di động Flutter đều tiêu thụ cùng một API REST và cổng WebSocket \texttt{/live}; nhật ký kiểm toán
được tách sang microservice gRPC riêng.}{Overall layered architecture
of OmniCast. Both the Next.js 15 web portal and the Flutter mobile application consume
the same REST API and \texttt{/live} WebSocket gateway; audit logging is separated into a dedicated gRPC
microservice.}}
\label{fig:architecture}
\end{figure*}

\bi{
Điểm thiết kế đáng chú ý là \emph{một tiến trình API duy nhất phục vụ
cả hai client}. Ranh giới giữa client web và client di động chỉ nằm ở
lớp vận chuyển, không nằm ở lớp nghiệp vụ. Nhờ vậy một quy tắc
nghiệp vụ chỉ được viết một lần: ví dụ, quy tắc ``hai chương trình
cùng kênh không được chồng lấn'' được kiểm tra trong
\texttt{ProgramsService.checkScheduleConflict} và có hiệu lực giống
nhau đối với cả thao tác tạo lịch từ trang quản trị web lẫn từ ứng
dụng di động.
}{
The notable design decision is a \emph{single API process serving both
clients}. The boundary between the web and mobile clients exists only
at the transport layer, never at the business layer. A business rule is
therefore written exactly once: for instance, the rule that "two
programmes on the same channel must not overlap" is enforced in
\texttt{ProgramsService.checkScheduleConflict} and behaves identically
whether the schedule is created from the web administration portal or
from the mobile application.
}

\subsection{\bi{Kiến trúc Truyền phát Trực tiếp (LiveStream) \& Bàn điều khiển Live Studio}{End-to-End LiveStream \& Live Studio Architecture}}

\bi{
Để đáp ứng yêu cầu phát sóng trực tiếp với độ trễ cực thấp và khả năng tương tác hai chiều giữa đạo diễn (\texttt{STAFF}/\texttt{ADMIN}) và khán giả (\texttt{VIEWER}), OmniCast thiết kế một kiến trúc truyền phát đa luồng chuyên biệt như minh hoạ tại Hình~\ref{fig:livestream_arch}.
}{
To satisfy ultra-low-latency live broadcasting and bi-directional interaction between studio producers (\texttt{STAFF}/\texttt{ADMIN}) and audiences (\texttt{VIEWER}), OmniCast implements a dedicated multi-pipeline live streaming architecture as shown in Fig.~\ref{fig:livestream_arch}.
}

\begin{figure*}[t]
\centering
\fitfull{\figLiveStreamArch}
\caption{\bi{Kiến trúc truyền phát trực tiếp (LiveStream) và bàn điều khiển Live Studio đầu-cuối: kết hợp 4 chế độ thu nhận luồng phát (Curated VOD, Camera/Screen WHIP WebRTC, Kho nội dung công khai, Điều hành tương tác) qua \texttt{LiveStudioService}, máy chủ MediaMTX và cổng Socket.IO \texttt{/live}.}{End-to-end LiveStream and Live Studio architecture: combining 4 studio ingestion modes (Curated VOD, Camera/Screen WHIP WebRTC, Verified Public Catalog, Host Control Console) through \texttt{LiveStudioService}, MediaMTX and the Socket.IO \texttt{/live} gateway.}}
\label{fig:livestream_arch}
\end{figure*}

\bi{
Kiến trúc LiveStream vận hành trên ba trụ cột kỹ thuật đồng bộ:
\begin{enumerate}[leftmargin=*,itemsep=2pt]
  \item \textbf{Bốn đường thu nhận luồng phát (Ingestion Pipelines) tại Live Studio:} Đạo diễn trên Web (\texttt{/studio/live}) hoặc Mobile (\texttt{LiveStudioScreen} \& \texttt{MobileStudioRecorderScreen}) có thể khởi tạo sóng trực tiếp qua 4 chế độ: (i) \emph{Curated Live} — chọn một hoặc nhiều bản ghi từ thư viện VOD nội bộ kèm bộ lọc chống vi phạm bản quyền (\emph{Anti-Piracy Guard}); (ii) \emph{Self-Record / Device Camera} — đẩy luồng trực tiếp từ ống kính máy quay phần cứng hoặc chia sẻ màn hình qua giao thức WebRTC WHIP tới MediaMTX (cổng \texttt{:8889/whip}) hoặc tải tệp ghi hình trực tiếp lên Cloudinary qua chữ ký số (\emph{Signed Direct Upload}); (iii) \emph{Verified Public Catalog} — nhập một chạm từ danh mục luồng HLS/MP4 công khai đã xác thực kèm kiểm tra \texttt{HEAD} probe; và (iv) \emph{Host Control Console} — gửi thông báo ghim, kiểm duyệt bình luận và chuyển chương trình đang phát tức thời.
  \item \textbf{Chuyển mã và Phân phối Đa giao thức (MediaMTX \& Direct Stream Engine):} Luồng WebRTC WHIP từ thiết bị quay được máy chủ MediaMTX chuyển đổi tức thời sang WHEP và HLS phân đoạn fMP4 độ trễ thấp (\texttt{:8888/<streamKey>/index.m3u8}). Ngoài ra, đối với các tệp lưu trên Google Drive, tầng backend và trình phát tự động phân giải đường dẫn sang \texttt{drive.usercontent.google.com/download?id=...\&confirm=t} để phát trực tiếp qua bộ đệm phần cứng mà không yêu cầu đăng nhập tài khoản Google.
  \item \textbf{Điều phối Phòng Live Thời gian thực qua Socket.IO (\texttt{/live}):} Cổng \texttt{LiveStreamGateway} quản lý các phòng theo sự kiện (\texttt{event:<eventId>}) và theo kênh (\texttt{channel:<channelId>}), đồng bộ tức thời số người đang xem thực tế (\texttt{live:viewerCount}), tin nhắn trò chuyện phân nhánh (\texttt{live:chat:new}), hiệu ứng biểu cảm nổi 0ms (\texttt{live:reaction:new}), thông báo ghim từ đạo diễn (\texttt{broadcaster:message}), quyên góp Super Chat (\texttt{live:donation:received}) và tín hiệu đổi luồng video tức thời (\texttt{live:program:changed}).
\end{enumerate}
}{
The LiveStream architecture operates across three synchronised technical pillars:
\begin{enumerate}[leftmargin=*,itemsep=2pt]
  \item \textbf{Four Studio Ingestion Pipelines:} Producers on Web (\texttt{/studio/live}) or Mobile (\texttt{LiveStudioScreen} \& \texttt{MobileStudioRecorderScreen}) can initiate broadcasts across four modes: (i) \emph{Curated Live} — selecting one or more recordings from the internal VOD library guarded by an Anti-Piracy filter; (ii) \emph{Self-Record / Device Camera} — pushing real-time camera or screen-share streams via WebRTC WHIP to MediaMTX (port \texttt{:8889/whip}) or uploading local recordings directly to Cloudinary using signed parameters; (iii) \emph{Verified Public Catalog} — one-click importing of verified public HLS/MP4 streams with server-side \texttt{HEAD} probing; and (iv) \emph{Host Control Console} — broadcasting pinned announcements, moderating chat and switching active programmes on the fly.
  \item \textbf{Multi-Protocol Transcoding \& Delivery (MediaMTX \& Direct Stream Engine):} Incoming WebRTC WHIP streams are remuxed by MediaMTX into WHEP and low-latency fMP4 HLS (\texttt{:8888/<streamKey>/index.m3u8}). Additionally, for Google Drive-hosted assets, the stream resolver rewrites share URLs to \texttt{drive.usercontent.google.com/download?id=...\&confirm=t}, enabling direct hardware-accelerated streaming without Google account prompts.
  \item \textbf{Real-Time Room Orchestration via Socket.IO (\texttt{/live}):} \texttt{LiveStreamGateway} manages per-event (\texttt{event:<eventId>}) and per-channel (\texttt{channel:<channelId>}) rooms, synchronising concurrent viewer counts (\texttt{live:viewerCount}), threaded chat (\texttt{live:chat:new}), 0ms floating emoji reactions (\texttt{live:reaction:new}), pinned host banners (\texttt{broadcaster:message}), Super Chat donations (\texttt{live:donation:received}) and zero-downtime stream switches (\texttt{live:program:changed}).
\end{enumerate}
}

\subsection{\bi{Cấu trúc module của NestJS}{NestJS Module Structure}}

\bi{
\texttt{AppModule} đăng ký 19 module nghiệp vụ cùng các module hạ tầng.
Mỗi module chịu trách nhiệm một miền nghiệp vụ và export các service
cần thiết; các controller chỉ gọi service và không chứa logic nghiệp
vụ. Cấu trúc thực tế của mã nguồn back-end được liệt kê ở
Bảng~\ref{tab:modules}.
}{
\texttt{AppModule} registers 19 business modules alongside the
infrastructure modules. Each module owns one business domain and
exports the services it needs; controllers only call services and hold
no business logic. Table~\ref{tab:modules} lists the actual
back-end source layout.
}

\begin{table}[t]
\caption{\bi{Cấu trúc module back-end NestJS.}{Back-end NestJS module
structure.}}
\label{tab:modules}
\centering
\footnotesize
\begin{tabular}{@{}lll@{}}
\toprule
\textbf{Module} & \bi{Trách nhiệm chính}{\textbf{Primary responsibility}} & \bi{Controller}{\textbf{Controller}} \\
\midrule
\texttt{auth}         & \bi{Đăng nhập, JWT, refresh, đăng nhập Google}{Login, JWT, refresh, Google sign-in} & \texttt{auth} \\
\texttt{channels}     & \bi{Kênh phát sóng, danh mục}{Broadcast channels, catalogue} & \texttt{channels} \\
\texttt{programs}     & \bi{Chương trình, EPG, Live Studio, Social}{Programmes, EPG, Live Studio, Social} & \texttt{programs} \\
\texttt{search}       & \bi{Tìm kiếm liên thực thể}{Cross-entity search} & \texttt{search} \\
\texttt{ai-curator}   & \bi{Biên tập AI, thống kê}{AI curation, statistics} & \texttt{ai-curator} \\
\texttt{audit-logger} & \bi{Nhật ký kiểm toán (gRPC)}{Audit logging (gRPC)} & \texttt{audit-logger} \\
\texttt{watchlist}    & \bi{Danh sách theo dõi}{Watchlist} & \texttt{watchlist} \\
\texttt{users}        & \bi{Quản lý người dùng}{User management} & \texttt{users} \\
\texttt{admin}        & \bi{Bảng điều khiển quản trị}{Admin dashboard} & \texttt{admin} \\
\texttt{billing}      & \bi{Gói dịch vụ, thanh toán}{Plans, payments} & \texttt{billing} \\
\texttt{notifications} & \bi{Thông báo, đẩy FCM}{Notifications, FCM push} & \texttt{notifications} \\
\texttt{newsletter}   & \bi{Đăng ký bản tin}{Newsletter subscription} & \texttt{newsletter} \\
\texttt{websocket}    & \bi{Cổng thời gian thực \texttt{/live}}{Realtime gateway \texttt{/live}} & --- \\
\texttt{cloudinary}   & \bi{Lưu trữ phương tiện}{Media storage} & \texttt{cloudinary} \\
\texttt{prisma}       & \bi{Kết nối cơ sở dữ liệu}{Database access} & --- \\
\texttt{health}       & \bi{Kiểm tra tình trạng}{Health probe} & \texttt{health} \\
\texttt{common}       & \bi{Bộ lọc, interceptor, Swagger}{Filters, interceptors, Swagger} & --- \\
\bottomrule
\end{tabular}
\end{table}

\subsection{\bi{Tích hợp giao thức và vận chuyển}{Protocol and Transport Integration}}

\bi{
Backend phục vụ bốn loại giao thức độc lập, mỗi loại phục vụ một nhóm
nhu cầu khác nhau:
}{
The back end serves four independent protocols, each serving a
distinct group of needs:
}

\begin{itemize}[leftmargin=*,itemsep=1pt]
  \item \bi{\textbf{REST/JSON qua HTTP/1.1.} Đường dẫn có tiền tố
  \texttt{/api/v1} và được tạo phiên bản bằng
  \texttt{VersioningType.URI}. Một middleware trong \texttt{main.ts}
  tự động chuyển hướng mọi yêu cầu dạng \texttt{/api/*} thành
  \texttt{/api/v1/*}, nhờ vậy các client cũ không gặp lỗi 404 khi
  phiên bản API được nâng cấp.}{
  \textbf{REST/JSON over HTTP/1.1.} Paths carry the \texttt{/api/v1}
  prefix and use \texttt{VersioningType.URI}. A middleware in
  \texttt{main.ts} automatically rewrites any \texttt{/api/*} request
  to \texttt{/api/v1/*}, so older clients never receive a 404 when the
  API version is upgraded.}

  \item \bi{\textbf{WebSocket (Socket.IO \texttt{/live}).} Cổng thời gian thực dùng Socket.IO với
  namespace \texttt{/live}, phục vụ đầy đủ các sự kiện điều hành phòng Live: theo dõi và bỏ theo dõi
  sự kiện/kênh (\texttt{subscribe:event}, \texttt{unsubscribe:event}, \texttt{subscribe:channel}, \texttt{unsubscribe:channel}), đếm mắt xem trực tiếp (\texttt{live:viewerCount}), chat trực tiếp (\texttt{live:chat} $\rightarrow$ \texttt{live:chat:new}), biểu cảm hạt nổi (\texttt{live:reaction} $\rightarrow$ \texttt{live:reaction:new}), thông báo ghim từ người phát (\texttt{broadcaster:message}),
  xoá tin nhắn vi phạm (\texttt{chat:delete} $\rightarrow$ \texttt{chat:deleted}), chuyển đổi chương trình tức thời (\texttt{live:program:changed}), quyên góp Super Chat (\texttt{live:donation:received}) và bắt tay WebRTC (\texttt{webrtc:offer}, \texttt{webrtc:answer}, \texttt{webrtc:ice-candidate}).}{
  \textbf{WebSocket (Socket.IO \texttt{/live}).} The realtime gateway uses Socket.IO on namespace
  \texttt{/live} and exposes full live room orchestration events: event and channel
  subscribe/unsubscribe (\texttt{subscribe:event}, \texttt{unsubscribe:event}, \texttt{subscribe:channel}, \texttt{unsubscribe:channel}), concurrent viewer tracking (\texttt{live:viewerCount}), live chat (\texttt{live:chat} $\rightarrow$ \texttt{live:chat:new}), floating reactions (\texttt{live:reaction} $\rightarrow$ \texttt{live:reaction:new}), broadcaster pinned messages (\texttt{broadcaster:message}),
  message moderation (\texttt{chat:delete} $\rightarrow$ \texttt{chat:deleted}), on-the-fly program switching (\texttt{live:program:changed}), Super Chat donations (\texttt{live:donation:received}) and the WebRTC handshake (\texttt{webrtc:offer}, \texttt{webrtc:answer}, \texttt{webrtc:ice-candidate}).}

  \item \bi{\textbf{gRPC.} Microservice ghi nhật ký kiểm toán nhận yêu
  cầu \texttt{LogAudit}, \texttt{GetLogsByUser},
  \texttt{GetLogsByAction}, \texttt{GetLogsByEntity} và
  \texttt{GetLogStats} qua \texttt{audit\_logger.proto}.}{
  \textbf{gRPC.} The audit-logging microservice receives
  \texttt{LogAudit}, \texttt{GetLogsByUser}, \texttt{GetLogsByAction},
  \texttt{GetLogsByEntity} and \texttt{GetLogStats} requests over
  \texttt{audit\_logger.proto}.}

  \item \bi{\textbf{WHIP/HLS/WebRTC.} MediaMTX nhận luồng phát qua
  WHIP và phân phối lại qua HLS (LL-HLS, fMP4) và WHEP. Cấu hình
  nằm trong \texttt{backend/mediamtx/mediamtx.yml}.}{
  \textbf{WHIP/HLS/WebRTC.} MediaMTX ingests over WHIP and redistributes
  over HLS (LL-HLS, fMP4) and WHEP, configured in
  \texttt{backend/mediamtx/mediamtx.yml}.}
\end{itemize}

\subsection{\bi{Kết nối cơ sở dữ liệu và giao thức ghi nhật ký}{Database Connection and Audit Logging}}

\bi{
Cơ sở dữ liệu là PostgreSQL do Supabase quản lý. Một chi tiết vận
hành đáng lưu ý: Supabase cung cấp hai cổng, cổng \texttt{5432} cho
kết nối trực tiếp và cổng \texttt{6543} cho transaction pooler. Hàm
\texttt{optimizeDatabaseUrl} trong \texttt{PrismaService} tự động đổi
cổng sang \texttt{6543} và đặt \texttt{connection\_limit} phù hợp khi
phát hiện pooler, tránh tình trạng cạn kết nối khi môi trường
serverless khởi động nhiều instance đồng thời.
}{
The database is Supabase-managed PostgreSQL. One operational detail is
worth noting: Supabase exposes two ports, \texttt{5432} for direct
connections and \texttt{6543} for the transaction pooler. The
\texttt{optimizeDatabaseUrl} helper in \texttt{PrismaService}
automatically switches to \texttt{6543} and sets a suitable
\texttt{connection\_limit} when it detects the pooler, avoiding
connection exhaustion when a serverless environment starts several
instances concurrently.
}

\begin{keybox}{\bi{Vì sao tách microservice ghi nhật ký kiểm toán?}{Why a
separate audit-logging microservice?}}
\bi{
Việc ghi nhật ký kiểm toán có hai đặc điểm không giống phần còn lại của
hệ thống: khối lượng ghi lớn và không chịu được sự chậm trễ. Nếu ghi
thẳng vào PostgreSQL, mỗi thao tác quản trị sẽ phải chờ một lần ghi
bổ sung. Tách sang microservice gRPC biến việc ghi nhật ký thành
thao tác bất đồng bộ theo kiểu fire-and-forget: \texttt{AuditLoggerService}
gọi RPC trong một \texttt{promise} không chặn, còn microservice tự ghi
vào cơ sở dữ liệu và duy trì bảng chỉ mục phục vụ truy vấn theo người
dùng, hành động, thực thể và thống kê tổng hợp.
}{
Audit logging has two properties the rest of the system does not
share: a high write volume and a tolerance for delay. Writing straight
into PostgreSQL would make every administrative operation wait on an
extra write. Moving it into a gRPC microservice turns auditing into a
non-blocking fire-and-forget operation: \texttt{AuditLoggerService}
issues the RPC inside a non-awaited promise, while the microservice
writes to the database and maintains index tables supporting queries by
user, action, entity and aggregate statistics.
}
\end{keybox}

\subsection{\bi{Cấu hình tích hợp ngoài}{External Integrations}}

\bi{
Bốn tích hợp bên ngoài được hỗ trợ, tất cả đều tự giáng cấp an toàn khi thiếu
khóa API: YouTube Data API v3 (lịch phát kênh), Twitch Helix (lịch phát
kênh), TMDB (làm giàu dữ liệu phim) và OpenAI (biên tập nội dung).
Cơ chế này cho phép toàn bộ hệ thống chạy đầy đủ chức năng trong môi
trường phát triển mà không cần bất kỳ khóa bên thứ ba nào.
}{
Four external integrations are supported, all of which degrade
gracefully when the API key is absent: YouTube Data API v3 (channel
schedules), Twitch Helix (channel schedules), TMDB (film metadata
enrichment) and OpenAI (content curation). This mechanism allows the
whole system to run at full capability in a development environment
without any third-party credential.
}
